%% title: Calibration split grouped by trap and night
%% target: Sections/Methods.tex
%% anchor: To rigorously capture calibration variance on finite distributions, we randomly partition the evaluation data into a calibration split (10\%) and a reporting split (90\%). Thresholds are calibrated strictly on the calibration split, while all reported metrics and model rankings are evaluated on the held-out reporting split across repeated Monte-Carlo trials.
%% action: after
%% why: A random split over trap crops puts frames of the same moth on both sides. Grouping removes that; measured, it moves calibrated scores by under a point and changes no conclusion.
%% source: journal/2026-10-22-rescoring-in-the-mini-metrics-convention.md, amendment of 2026-10-23 (65 probe runs: +0.8 on average, range -0.5 to +2.3); dev/093_rescore_mini_metrics.py --grouped
%% check: Asger's own trap numbers need the same grouped split before they are compared with ours.
%% status: ready
On the trap images the split is made by capture group: a random 10\,\% of the (trap, night) groups calibrates the thresholds and the rest is scored, so that no moth photographed on one night appears on both sides.
On the probe set this changes calibrated macro-$\text{F}_1$ by 0.8 points on average (between $-0.5$ and $+2.3$ over 65 models) and no comparison between models.
